\begin{figure}
  \centering
  \begin{tikzpicture}[node distance=2cm]
    \draw node at (0,0)[name=insieme]{set};
    \draw node at (-1,-1)[name=sv]{vector space};
    \draw node at (-1,-3)[name=sn]{normed space};
    \draw node at (-1,-5)[name=euclideo]{Euclidean space};
    \draw node at (2,-1)[name=st]{topological space};
    \draw node at (2,-2)[name=sm]{metric space};
    \draw node at (2,-3.5)[name=completo]{complete space};
    \draw node at (1,-4.5)[name=banach]{Banach};
    \draw node at (1,-6)[name=hilbert]{Hilbert};
    \draw[->](sv) -- (insieme);
    \draw[->](st) -- (insieme);
    \draw[->](sn) -- (sv);
    \draw[->](sm) -- (st);
    \draw[->](euclideo) -- (sn);
    \draw[->](hilbert) -- (euclideo);
    \draw[->](hilbert) -- (banach);
    \draw[->](banach) -- (completo);
    \draw[->](completo) -- (sm);
    \draw[->](sn) -- (sm);
    \draw[->](banach) -- (sn);

    \begin{scope}[black!40]
    \draw node at (-2,0)[name=insiemeT]{$\in,\subset,\cup$};
    \draw (insiemeT) -- (insieme);
    \draw node at (-3,-1)[name=svT]{$+,\cdot$};
    \draw (svT) -- (sv);
    \draw node at (-3,-3)[name=snT]{$\abs{v}$};
    \draw (snT) -- (sn);
    \draw node at (-3,---5)[name=euclideoT]{$\langle v,w\rangle$};
    \draw (euclideoT) -- (euclideo);
    \draw node at (5,0)[name=stT]{$\lim, \stackrel\circ A, \bar A$};
    \draw (stT) -- (st);
    \draw node at (5,-1.75)[name=smT]{$d(x,y)$, Lip};
    \draw (smT) -- (sm);
    \end{scope}

    \begin{scope}[black!75]
    \draw node at (5,-1)[name=Rbar,left]{$\bar \RR$};
    \draw[->](Rbar) -- (st);
      \draw node at (5,-2.5)[name=Q,left]{$\QQ$};
      \draw[->](Q) -- (sm);
      \draw node at (5,-4)[name=Sn,left]{$\mathbb S^n$};
      \draw[->](Sn) -- (completo);
      \draw node at (4.5,-5)[name=Cn,left]{$\mathcal C^n$, $L^p$, $\ell^p$};
      \draw[->](Cn) -- (banach);
      \draw node at (4,-6.5)[name=L2,left]{$L^2$, $\ell^2$};
      \draw[->](L2) -- (hilbert);
      \draw node at (1,-7)[name=Rn]{$\RR^n$};
      \draw[->](Rn) -- (hilbert);
    \end{scope}
  \end{tikzpicture}
\caption{Mathematical structures abstracting the space $\RR^n$.
The black arrow means: ``is a'',
the dark arrow: ``is an example of'',
the light line: ``is a typical operation of''.}
\label{fig:spazi}
\end{figure}

Abstraction is a typical activity in mathematics. The results we have seen so
far mostly concern the space $\RR$ of real numbers and partly the space $\CC$ of
complex numbers. However, it is evident that many of these results extend to
other more general situations. Identifying the fundamental properties underlying
a theorem is perhaps the most important activity performed by mathematics.
Typically, identifying these properties not only generalizes a result but also
simplifies its proof. Indeed, once the minimal hypotheses of the theorem are
identified, the tools at our disposal are also reduced, making it easier to
understand which ones should be used. On the other hand, we cannot take this
attempt at generalization to the extreme, as it quickly becomes apparent that
the array of spaces that gradually cover all possible types becomes enormous,
and keeping all the definitions in mind becomes too burdensome.

Refer to Figure~\ref{fig:spazi} for a relational overview of the spaces we will
now introduce. Vector spaces (linear algebra) are assumed to be known, as they
are usually covered in geometry courses. Topological spaces will only be briefly
mentioned as they will not be directly useful for our purposes, and a complete
treatment would require a significant amount of time. For the same reason, other
very relevant spaces (such as groups and differential manifolds) are completely
absent from this classification.

\begin{definition}[metric space]
\mymark{**}
\label{def:distanza}
We will say that $d\colon X\times X\to\RR$ is a \emph{distance}%
\mymargin{distance}%
\index{distance} on $X$
if for every $x,y,z\in X$ the following properties hold:
\begin{enumerate}
\item
  $d(x,y)\ge 0$ (positivity);
\item
  $d(x,y)\le d(x,z) + d(z,y)$ (triangle inequality);
\item
  $d(x,y)=0$ if and only if $x=y$ (separation);
\item
  $d(x,y) = d(y,x)$ (symmetry).
\end{enumerate}

If $d$ is a distance, we will say that $X$ is a
\emph{metric space}
\mymargin{metric space}%
\index{space!metric}%
(with distance $d$).
\end{definition}

Observe that from the triangle inequalities:
\[
  d(x,z) \le d(x,y) + d(y,z), \qquad
  d(y,z) \le d(x,y) + d(x,z)
\]
one obtains the \emph{reverse triangle inequality}:
\mymargin{reverse triangle inequality}
\index{inequality!triangle!reverse}
\[
  d(x,y) \ge \abs{d(x,z) - d(y,z)}.
\]

\begin{definition}[convergence]%
  \label{def:convergenza_metrica}%
\mymark{*}%
If $x_n \in X$ is a sequence of points in a metric space $X$
and $x\in X$, we will say that $x_n$ converges to $x$ and write
\[
  x_n \to x
\]
if $d(x_n,x)\to 0$ as $n\to +\infty$.
\end{definition}

Note that the usual convergence in $\RR$ is nothing other than the convergence
of $\RR$ viewed as a metric space with the Euclidean distance $d(x,y) =
\abs{x-y}$.

With the artifice $d(x_n,y_n) - d(x,y) = d(x_n,y_n) - d(x_n,y) + d(x_n,y) -
d(x,y)$ and using the reverse triangle inequality, one can easily solve the
following.

\begin{exercise}[continuity of distance]
Prove that if $x_n\to x$ and $y_n\to y$, then $d(x_n,y_n) \to d(x,y)$.
\end{exercise}

If we wish to define a distance on a vector space $V$, it will be natural to
require that the distance is compatible with the vector space structure, and
thus that it is invariant under translations and homogeneous with respect to
dilations. With this requirement, the distance $d(x,y)$ between two points must
be equal to the distance of $x-y$ from the origin. The distance from the origin
is called a \emph{norm} and is the basis of the following definition.

\begin{definition}[normed space]
\mymark{*}
\label{def:norma}
Let $V$ be a vector space over the field $\RR$.
A function $\phi \colon V\to \RR$ is said to be a \emph{norm}%
\mymargin{norm}%
\index{norm} on $V$ if
for every $v,w\in V$ and for every $\lambda \in \RR$
the following properties hold:
\begin{enumerate}
\item
  $\phi(v) \ge 0$ (positivity);
\item
  $\phi(\lambda v) = \abs{\lambda} \cdot \phi(v)$ (homogeneity and symmetry);
\item
  $\phi(v+w) \le \phi(v) + \phi(w)$
  (triangle inequality);
\item
  $\phi(v)=0$ if and only if $v=0$ (separation).
\end{enumerate}

If $\phi$ is a norm on $V$, we will say that
$V$
is a \emph{normed vector space}%
\mymargin{normed}%
\index{normed} by $\phi$.

If $\phi$ is a norm, the function $d(v,w) = \phi(v-w)$
is clearly a distance called the
\emph{induced distance}
\mymargin{induced distance}%
\index{distance!induced}%
by $\phi$.
In particular, every normed space is also a metric space with respect to the
distance induced by the norm.
%Moreover, the triangle inequality
%ensures that the norm $\phi\colon V\to \RR$ is a continuous function.
\end{definition}

Usually, the notations $\abs{v}$ or $\Abs{v}$ are used to indicate
a norm $\phi(v)$.

Since a normed space is also a metric space, convergence is defined in normed
spaces. It is easy to verify that the norm is continuous with respect to such
convergence, in the sense that if $v_k \to v$
(that is, $\phi(v_k-v)\to 0$) then $\phi(v_k)\to \phi(v)$.

\begin{definition}[Euclidean space]
\label{def:prodotto_scalare}%
Let $V$ be a vector space over the field $\RR$.
A function $b\colon V\times V \to \RR$ is said to be a \emph{positive definite
inner product}%
\mymargin{positive definite inner product}%
\index{product!inner}%
\index{definite!positive}
on $V$ if
$b$ is a bilinear, symmetric, and positive definite form, that is, if
for every $u,v,w\in V$ and for every $\lambda, \mu \in \RR$
the following properties hold:
\begin{enumerate}
\item $b(v,v) \ge 0$ (positivity);
\item $b(\lambda u + \mu v, w) = \lambda b(u,w) + \mu b(v,w)$
and $b(u,\lambda v+ \mu w) = \lambda b(u,v) + \mu b(u,w)$
(bilinearity);
\item $b(v,w) = b(w,v)$ (symmetry);
\item $b(v,v) = 0$ if and only if $v=0$ (non-degeneracy).
\end{enumerate}

If $b$ is a positive definite inner product on $V$, we will say that $V$ is a
\emph{Euclidean space} (with inner product $b$).

Usually, the notations $v\cdot w$, $(v,w)$, $\langle v,w\rangle$
or $\langle v \vert w\rangle$ are used to denote an inner product $b(v,w)$.
\end{definition}

If $b$ is a positive definite inner product on $V$, the following theorem
guarantees that
the function $\phi(v) = \sqrt{\langle v,v\rangle}$ is a norm on $V$.

Thus, a Euclidean space naturally has the structure of a normed space and a
metric space.

\begin{theorem}[properties of the inner product]
\label{th:spazio_euclideo}%
Let $V$ be a Euclidean space with a (positive definite) inner product $\langle
v, w\rangle$.
Define $\abs{v} = \sqrt{\langle v,v\rangle}$.
Then $\abs{v}$ is a norm on $V$
and for every $v,w\in V$ the following properties hold:
\begin{enumerate}
\item \emph{binomial expansion}%
\mymargin{binomial expansion}%
\index{expansion!binomial}
\begin{equation}\label{eq:binomio_vettoriale}%
  \abs{v+w}^2 = \abs{v}^2 + 2 \langle v,w\rangle + \abs{w}^2;
\end{equation}
\item \emph{Pythagorean theorem}%
\mymargin{Pythagorean theorem}%
\index{theorem!Pythagorean}
\index{Pythagoras!theorem}%
\begin{equation}\label{eq:Pitagora}
  \langle v,w\rangle=0 \implies \abs{v-w}^2 = \abs{v}^2 + \abs{w}^2;
\end{equation}
\item \emph{Young's inequality}%
\mymargin{Young's inequality}%
\index{inequality!Young's}
\index{Young!inequality of}%
\begin{equation}\label{eq:Young}
  \langle v,w\rangle \le \frac{\abs{v}^2 + \abs{w}^2}{2};
\end{equation}
\item \emph{Cauchy-Schwarz inequality}%
\mymargin{Cauchy-Schwarz inequality}%
\index{inequality!Cauchy-Schwarz}%
\index{Cauchy-Schwarz!inequality of}%
\begin{equation}\label{eq:Cauchy_Schwarz}
   \langle v,w \rangle \le \abs{v} \cdot \abs{w};
\end{equation}
\item \emph{parallelogram law}%
\mymargin{parallelogram law}%
\index{law!parallelogram}
\index{parallelogram!law of}
\begin{equation}\label{eq:parallelogramma}
  \abs{v+w}^2 + \abs{v-w}^2 = 2 \abs{v}^2 + 2\abs{w}^2;
\end{equation}
\item \emph{continuity}
\[
  \text{if $v_k\to v$ and $w_k\to w$, then
  $\langle v_k,w_k\rangle \to \langle v,w \rangle$}.
\]
\end{enumerate}
\end{theorem}
%
\begin{proof}
First, observe that $\abs{v}$ is well-defined for every $v\in V$
since the inner product $\langle v,v\rangle$ is by definition never negative.
Then, in general, using bilinearity, we have the usual expansion of the square
of the binomial:
\begin{align*}
  \abs{v+w}^2
  &= \langle v+w, v+w\rangle
  = \langle v+w, v\rangle + \langle v+w, w\rangle \\
  &= \langle v,v\rangle + 2\langle v,w \rangle + \langle w,w\rangle
  = \abs{v}^2 + 2\langle v,w \rangle + \abs{w}^2.
\end{align*}
The Pythagorean theorem follows immediately.
But then Young's inequality is easily obtained:
\[
  0 \le \abs{v-w}^2 = \abs{v}^2 + \abs{w}^2 - 2 \langle v,w\rangle
\]
and the parallelogram law:
\[
  \abs{v+w}^2 + \abs{v-w}^2
  = \abs{v}^2 + 2\langle v,w \rangle + \abs {w}^2
    + \abs{v}^2 - 2\langle v,w\rangle + \abs{w}^2.
\]
Now, if $\abs{v}=\abs{w}=1$, Young's inequality tells us that
\[
  \langle v,w \rangle \le 1.
\]
But then for every $v\neq 0$ and $w \neq 0$, the Cauchy-Schwarz inequality is
obtained:
\[
\frac{\langle v,w\rangle}{\abs{v}\cdot \abs{w}}
 = \left\langle \frac{v}{\abs{v}} , \frac{w}{\abs{w}}\right\rangle
 \le 1.
\]
If instead $v=0$ or $w=0$, the Cauchy-Schwarz inequality is obvious since
for every $u\in V$, we have
$\langle 0, u\rangle = \langle u,0\rangle = 0$ by bilinearity.

The properties of positivity, homogeneity, symmetry, and separation of 
the norm $\abs{v}$ follow directly from the analogous properties 
of the inner product $\langle v,w\rangle$.  
The triangle inequality follows from the 
Cauchy-Schwarz inequality,
indeed:
\[
 \abs{v+w}^2 = \langle v+w, v+w\rangle 
 = \abs{v}^2 + \abs{w}^2 + 2 \langle v,w\rangle 
 \le \abs v^2 + \abs w^2 + 2 \abs{v}\cdot \abs{w}
 = \enclose{\abs{v} + \abs{w}}^2
\]
from which $\abs{v+w} \le \abs{v} + \abs{w}$.


Regarding continuity,
remember that $v_k \to v$ means $\abs{v_k-v}\to 0$.
We can then write
\begin{align*}
  \langle v_k,w_k\rangle - \langle v,w\rangle
&= \langle v_k,w_k\rangle - \langle v, w_k\rangle +  \langle v,w_k\rangle -
\langle v,w\rangle \\
&= \langle v_k - v, w_k\rangle + \langle v, w_k -w\rangle
\end{align*}
and, using Cauchy-Schwarz, if $v_k\to v$ and $w_k\to w$, we obtain
\begin{align*}
\abs{\langle v_k,w_k\rangle - \langle v,w\rangle}
&\le \abs{v_k-v}\cdot \abs{w_k} + \abs{v}\cdot \abs{w_k-w}\\
&\to 0\cdot \abs{w} + \abs{v}\cdot 0 = 0.
\end{align*}
\end{proof}

The vector space $\RR^n$ has a canonical Euclidean structure,
as in the following.

\begin{definition}[Euclidean structure of $\RR^n$]%
\label{def:124124}%
\mymark{*}%
A vector $\vec x\in \RR^n$ is defined as an $n$-tuple of real numbers:
\[
  \vec x = (x_1, \dots, x_n).
\]
On $\RR^n$, we can then define the inner product:
\index{inner product!of $\RR^n$}
\[
  \vec x\cdot \vec y
   = \sum_{k=1}^n x_k y_k
   = x_1 y_1 + \dots + x_n y_n
\]
This inner product induces the Euclidean norm:
\index{norm!Euclidean of $\RR^n$}%
\[
  \abs{\vec x} = \sqrt{\vec x \cdot \vec x}
  = \sqrt{x_1^2 + x_2^2+ \dots + x_n^2}.
\]
The distance induced by this norm is called the \emph{Euclidean distance}%
\mymargin{Euclidean distance}%
\index{distance!Euclidean}:
\[
  d(\vec x, \vec y)
  = \abs{\vec x-\vec y}
  = \sqrt{(x_1-y_1)^2 + \dots + (x_n-y_n)^2}.
\]

In the case $n=1$, the norm coincides with the absolute value, which
justifies the use of the same notation.

If we identify $\CC$ with $\RR^2$ by associating the complex number $z=x+iy$
with the point $(x,y)\in \RR^2$, we can observe that the Euclidean norm
coincides with the modulus of the complex number:
\[
  \abs{z} = \sqrt{x^2 + y^2} = \abs{(x,y)}.
\]

Thus, $\RR$, $\CC$, $\RR^n$,
are Euclidean spaces, normed spaces, and metric spaces with respect to the
canonical
Euclidean structure.
\end{definition}

\begin{example}[Manhattan distance]
On $\RR^2$, we can define a norm, called the \emph{Manhattan norm}%
\mymargin{Manhattan norm}%
\index{norm!Manhattan},
\index{Manhattan!norm}
as follows:
\[
  \phi(\vec x) = \abs{x_1} + \abs{x_2}.
\]
The induced distance $d(p,q)$
\index{distance!Manhattan}%
\index{Manhattan!distance}%
represents the length of the shortest path to
go from $p$ to $q$ moving only horizontally or vertically
(as if we were on the streets of Manhattan).

The Manhattan norm is not Euclidean in the sense that it is not possible to
define
an inner product that induces this norm. Indeed, if such an
inner product existed, the parallelogram law
\eqref{eq:parallelogramma} would have to hold, but we observe that choosing $v=(1,0)$ and $w=(0,1)$
we have
\[
  \phi(v+w)^2 + \phi(v-w)^2
  = 8
  \neq 4
  = 2\phi(v)^2 + 2\phi(w)^2.
\]

The set $B_R(\vec x_0) = \ENCLOSE{\vec x \in \RR^2\colon \phi(\vec x-\vec x_0) <
R}$
of points in $\RR^2$ that are less than $R$ away from the point $\vec x_0$
(see definition~\ref{def:palla})
is a square
with diagonals, of length $2R$, parallel to the coordinate axes.
If as a norm $\phi$ we had chosen the canonical Euclidean norm of $\RR^2$,
the set $B_R(\vec x_0)$ would have been a circle of radius $R$
centered at $\vec x_0$.
In general, norms induced by an inner product are recognized
by the shape of these sets: only when ellipses
(or ellipsoids if we are in higher dimension) are obtained, the norm is
Euclidean.
For other norms, one may obtain generic convex symmetric sets
with respect to the center $\vec x_0$.
\end{example}

\begin{example}[$p$-norm]
\label{ex:norma_p}%
For every $p\ge 1$, one can define on $\RR^n$ the norm
\mymargin{$p$-norm}%
\index{$p$-norm}%
\index{norm!$p$ in $\RR^n$}%
\[
  \abs{\vec x}_p = \sqrt[p]{\abs{x_1}^p + \abs{x_2}^p + \dots + \abs{x_n}^p}.
\]

One can also define
\[
  \abs{\vec x}_\infty = \lim_{p\to +\infty} \abs{\vec x}_p
   = \max\ENCLOSE{\abs{x_1}, \abs{x_2}, \dots, \abs{x_n}}.
\]

For $p=2$, one obtains the Euclidean norm of definition~\ref{def:124124}
$\abs{v}_2=\abs{v}$.
For $p=1$
on $\RR^2$, one obtains the Manhattan norm.
For $p=+\infty$, one obtains again
the Manhattan norm up to a 45-degree rotation and a scaling
by a factor of $\sqrt 2$:
\[
  \abs{(x_1,x_2)}_\infty
  = \abs{\enclose{\frac{x_1-x_2} 2, \frac {x_1+x_2}2}}_1
\]

It could be shown that only for $p=2$ is the norm $\abs{\vec x}_p$ induced
by an inner product, as only for $p=2$ is the
parallelogram law~\eqref{eq:parallelogramma} valid.
\end{example}